Los circuitos entre neuronas biológicas y modelos computacionales requieren sinapsis artificiales cuyos parámetros, no reutilizables entre experimentos, reproduzcan ciertas dinámicas. La parametrización del modelo de sinapsis química graduada de Golowasch et al. (que separa la corriente en componentes lenta y rápida) es compleja debido a su elevado número de parámetros, el acoplamiento entre ellos y las restricciones temporales de los experimentos con neuronas vivas. Este trabajo aborda este problema mediante un algoritmo de parametrización capaz de inducir acoplamientos en fase o antifase, garantizando corrientes sinápticas realistas.

Tras revisar en el estado del arte los modelos sinápticos y algoritmos de optimización, se elige la optimización bayesiana por su eficiencia muestral y su capacidad para tratar el solapamiento paramétrico. Luego, se explica la función de evaluación, basada en la descomposición frecuencial del potencial presináptico, separándolo en una componente lenta y otra rápida (\textit{spikes}), lo que permite evaluar independientemente cada componente de la corriente combinando similitud de forma (con este potencial) y adecuación de extremos. Más tarde, al comentar la optimización bayesiana, se explica que emplea un proceso gaussiano con \textit{kernel Squared Exponential} y \textit{Automatic Relevance Determination}. Además, se incorporan en los rangos de exploración reparametrizaciones y escalas logarítmicas para mitigar el solapamiento, y límites dinámicos para adaptar la exploración a cada experimento.

Más adelante se comentan las implementaciones desarrolladas: una \textit{offline} que opera con registros neuronales de forma unidireccional y luego se evaluará comunicando un registro biológico con un modelo de Hindmarsh-Rose; y un módulo \textit{online} integrado en \textit{RTXI} para comunicación bidireccional (que se diseña para no afectar el funcionamiento del tiempo real), cuya evaluación se hará entre dos modelos. Se analizará, tanto buscando fase como antifase, la convergencia del algoritmo, la evolución de las métricas, la importancia de los parámetros aprendida y las dinámicas neuronales resultantes.

Los resultados demuestran una convergencia adecuada, induciendo los acoplamientos deseados mediante corrientes realistas. Además de la versión \textit{offline}, que termina en segundos, la \textit{online} lo hace en menos de 16 minutos, duración compatible con los experimentos biológicos.